%-----------------------------------------------------------
%   Capítulo 3 - Cenarios de Aplicaçao e Analise de Requisitos
%-----------------------------------------------------------
\chapter{Análise de Requisitos}


\section{Application Scenarios}

The application scenarios for the solar park inspection system using a wheeled mobile robot include the following:

\begin{itemize}
    \item Cable inspection: Checking the condition and integrity of electrical cables in the solar park.
    \item Vegetation monitoring: Observing plant growth that may interfere with solar panels.
    \item Site surveillance: General monitoring of security and intruder presence in the park.
    \item Thermal surveillance: Detection of thermal anomalies such as hot spots in solar panels through thermal images.
\end{itemize}

\section{Requirements}

The requirements analysis for the system development considers several fundamental aspects:

\begin{itemize}
    \item \textbf{Wheeled mobile robot}: Use of a wheeled robot suitable for navigation on uneven solar park terrains.
    \item \textbf{Sensors}: 
        \begin{itemize}
            \item For navigation: RTK-GPS, IMU (Inertial Measurement Unit), odometry.
            \item For localization: LiDAR, stereo cameras, ultrasonic sensors.
            \item For inspection: Visual cameras, thermal cameras, proximity sensors.
        \end{itemize}
    \item \textbf{Framework}: Use of ROS (Robot Operating System) for sensor integration, control, and data processing.
    \item \textbf{Predefined information}: Map of the solar park to apply geo-referencing and path planning.
\end{itemize}